% ============================================================================
% SchizoStudy v5.0 — сжатый технический отчёт (русский, без рисунков).
% Компилируется xelatex (как brochure.tex). Шрифты лежат рядом, в fonts/.
% ============================================================================
\documentclass[11pt,a4paper]{article}
\usepackage{fontspec}
\usepackage{polyglossia}
\setdefaultlanguage{russian}\setotherlanguage{english}
% Шрифты подключаются по имени файла, устанавливать в систему их не нужно.
\newcommand{\optserif}{Path=fonts/,Extension=.ttf,UprightFont=*-Regular,BoldFont=*-Bold,ItalicFont=*-Italic,BoldItalicFont=*-BoldItalic,Ligatures=TeX}
\newcommand{\optmono}{Path=fonts/,Extension=.ttf,UprightFont=*-Regular,BoldFont=*-Bold,ItalicFont=*-Italic,BoldItalicFont=*-BoldItalic,Scale=0.9,Ligatures=TeX}
\setmainfont{LiberationSerif}[\optserif]\newfontfamily\cyrillicfont{LiberationSerif}[\optserif]
\setmonofont{LiberationMono}[\optmono]\newfontfamily\cyrillicfonttt{LiberationMono}[\optmono]
\usepackage[a4paper,margin=1.9cm]{geometry}
\usepackage{enumitem,amsmath,amssymb}
\usepackage{booktabs,titlesec,amsthm,array}
\usepackage[hidelinks]{hyperref}
\setlist[itemize]{leftmargin=1.0em,itemsep=0pt,topsep=1pt}
\setlist[enumerate]{leftmargin=1.2em,itemsep=0pt,topsep=1pt}
\setlength{\parskip}{1pt}
\emergencystretch=2em
\titleformat{\section}{\large\bfseries}{\thesection.}{0.4em}{}
\titlespacing*{\section}{0pt}{7pt plus 1pt}{3pt}
\titleformat{\subsection}{\normalsize\bfseries}{\thesubsection.}{0.4em}{}
\titlespacing*{\subsection}{0pt}{5pt}{2pt}
\newtheorem{proposition}{Предложение}[section]
\newtheorem{definition}[proposition]{Определение}
\theoremstyle{remark}
\newtheorem{remark}[proposition]{Замечание}
\newcommand{\code}[1]{\texttt{\detokenize{#1}}}

\begin{document}

\begin{center}
{\large\bfseries Отчёт}\\[2pt]
{\small Предзарегистрированная \emph{in silico} программа опровержения
модели повреждения парвальбуминовых интернейронов на основе редокс-переключателя}\\[2pt]
{\small Никита Филиппов \textperiodcentered\
\href{https://github.com/stockman2/SchizoStudy-v5}{github.com/stockman2/SchizoStudy-v5}
\textperiodcentered\ 24 октября 2026 г.}
\end{center}
\vspace{-4pt}

\noindent\textbf{Аннотация.} Построена механистическая ОДУ-модель окислительного
повреждения парвальбуминовых (PV) интернейронов: бистабильный редокс-переключатель
«пероксиредоксин--сульфиредоксин», двухступенчатая индуцируемая защита
($\tau_m=12$\,ч, $\tau_s=63$\,ч) с праймингом, асимметрией и поклеточной ёмкостью,
переключатель сна и бодрствования Филлипса--Робинсона и кора Вильсона--Коуэна,
замыкающая петлю через растормаживание, --- после чего большая часть усилий потрачена
на попытки её опровергнуть. Каждая гипотеза записана как нумерованное правило; каждый
порог выведен из ожиданий и дисперсий до счёта; вердикты лежат в журнале, который
только дополняется, и
никогда не пересматриваются. Из 128 предзарегистрированных частей в 26 правилах 66
выполнены, 48 не выполнены, 13 непроверяемы, одна ниже собственной аналитической
вилки. Модель описывает возникновение и окончание эпизода под окислительной
нагрузкой; главный результат --- её неудачи.

\section{Объект изучения}\label{sec:intro}
\emph{Гипотеза.} PV-интернейроны удерживают возбуждение коры в рамках, задают
гамма-ритм и платят за это самым высоким расходом энергии в сети --- и, значит, самым
большим потоком перекиси. Часть антиоксидантной защиты клетка производит, только пока
работает её NMDA-рецептор --- датчик активности; слабый рецептор означает меньше защиты
и больше окислителей, а окислители подавляют сам рецептор: положительная петля.
Перегруженная PV-клетка не гибнет, а \emph{замолкает} (меньше парвальбумина и GAD67),
кора растормаживается, нагрузка растёт: вторая петля; расторможённый гиппокамп повышает
чувствительность дофаминовой системы (схема Грейса) --- мост к позитивным симптомам.
Каждая стрелка подтверждена опытами по отдельности; вся цепочка целиком --- никогда.
Количественный вопрос: достаточно ли усиление петли, чтобы она \emph{замкнулась}, то
есть сделала клетку бистабильной? Бистабильность --- ровно то, что превращает медленное
накопление в резкий переход с гистерезисом, то есть в клиническую форму течения
(долгий продром, резкий первый эпизод, эпизоды длиной от часов до суток).

\emph{Стратегия.} Собрать минимальную модель; затем систематически её опровергать:
нумерованные правила с частями, пороги --- только из аналитических ожиданий, журнал
вердиктов, который только дополняется, побитово сверенный код, замороженный эталон. Версии: v3 добавила
переключатель сна, кору и популяцию; v4 --- двухступенчатую защиту, прайминг, лестницу и
статистическую дисциплину; v5 --- асимметричную защиту, три фазы популяции, слитое ядро
и два правила о петлях (Ф26--Ф29).

\section{Модель}\label{sec:model}
Модель --- стопка из четырёх слоёв с разделёнными масштабами времени: клетка
(минуты--сутки), переключатель сна (часы), кора (адиабатически исключена), популяция (через один
скаляр).

\subsection{Редокс-ядро}\label{sec:core}
Пусть $x\in[0,1]$ --- доля рабочего пероксиредоксина, $c$ --- защита, поддерживаемая
рецептором, $u$ --- перекись, $\nu$ --- нагрузка. Перекись быстрая (алгебраический
баланс); динамика такова:
\begin{equation}\label{eq:core}
\begin{aligned}
&\nu=\varepsilon u+\frac{xu}{1+u},\qquad
\varepsilon=e_0c^{\gamma},\qquad
N=\frac{N_0}{1+u/K_N},\\
&x'=c^{m}(1-x)-\mu_0\frac{xu^2}{1+u},\qquad
c'=\frac{1}{T}\Bigl(N+\frac{\beta u}{K_B+u}-c\Bigr).
\end{aligned}
\end{equation}
Первая строка и уравнение для $x$ --- измеренная кинетика (Сельваджо и соавт., 2018);
уравнение для $c$ --- допущение о форме. $T$ --- время жизни защитных белков в единицах
времени починки; $N_0<1$ --- врождённо слабый рецептор.

\begin{proposition}[Бистабильность]
Переокисленный пероксиредоксин чинится медленно, поэтому доля рабочего фермента с
ростом перекиси \emph{убывает}, а очистка \emph{требует} фермента: обратная связь
в~\eqref{eq:core} даёт две седло-узловые складки и окно гистерезиса по нагрузке, внутри
которого «защита включена» и «защита выключена» устойчивы одновременно. Измеренные
клетки лежат вне окна; чтобы попасть внутрь, починку нужно ослабить в $1{,}6$--$22$
раза --- ровно то состояние PV-клеток при слабом NMDA-сигнале, которое предполагает
гипотеза.
\end{proposition}

\emph{Наборы параметров.} Девять параметров разыгрываются ($e_0$, $\mu_0$, $K_N$,
$K_B$, $T$ --- равномерно по логарифму; $m\in[1,2]$, $\gamma\in[0,1]$,
$N_0\in[0{,}3;1]$ --- равномерно), причём $\beta=N_0\cdot U(0,3)$ (измеренная индукция
до $2{,}85$ раза). Набор принимается, только если он бистабилен: отношение границ окна
больше $1{,}05$ и два устойчивых состояния при среднегеометрической нагрузке. Принятые
наборы --- это \emph{потенциальные больные} модели: только тот набор, чья клетка уже
сидит в бистабильном режиме (ослабленная защита --- то самое состояние, которое
гипотеза постулирует для PV-клеток при слабой NMDA-передаче), вообще может получить
описываемое моделью повреждение; остальные наборы здоровы по построению и не
изучаются. Из
$K=200\,000$ кандидатов отбирается $n_{\mathrm{sel}}=2400$ принятых наборов
(детерминированно по сиду, с кэшем); средняя нагрузка стоит на $q=0{,}10$ окна.

\subsection{Индуцируемая защита}\label{sec:defence}
\emph{Две ступени.} Сначала сигнал индукции $m$, затем белок $s$, прибавляемый к $c$:
\begin{equation}\label{eq:twostage}
\dot m=\frac{1}{\tau_m}\Bigl(\beta\frac{u}{K_B+u}-m\Bigr),\qquad
\dot s=\frac{1}{\tau_s}(m-s),\qquad
\tau_m=12\,\mathrm{ч},\ \tau_s=63\,\mathrm{ч},
\end{equation}
с откликом на ступень
$s(t)=1-(\tau_s e^{-t/\tau_s}-\tau_m e^{-t/\tau_m})/(\tau_s-\tau_m)$ --- всегда с
перегибом; половина уровня набирается к 57\,ч. Стационар (а с ним окно, складки и закон
разрывов) не меняется; меняется только ход во времени.

\emph{Прайминг.} След $p$, набираемый на том же сигнале, что движет индукцию, усиливает её:
$\dot p=(u/(K_B+u)-p)/\tau_p$, $\beta_{\mathrm{eff}}=\beta(1+Ap)$, $A=2$ (потолок
$\times3$ измерен; $\tau_p$ не измерена и перебирается).

\emph{Асимметрия.} Постоянные подъёма и спада $s$ разведены,
$r=\tau_{\downarrow}/\tau_{\uparrow}$.

\begin{proposition}[Асимметричный интегратор]\label{prop:wr}
Для двухуровневого суточного входа, который высок в течение доли $f$ суток,
стационарное среднее переменной-интегратора стоит в точке
\begin{equation}\label{eq:wr}
w(r)=\frac{fr}{fr+1-f},
\end{equation}
$f=0{,}705$ (измеренная доля бодрствования). Асимметрия превращает интегратор в
детектор пика ($r>1$) или детектор дна ($r<1$); в~\eqref{eq:wr} нет ни одного
подогнанного числа.
\end{proposition}

\emph{Петля ёмкости.} Поклеточно: пока клетка повреждена, $N_0\to N_0(1-a_n)$ (потеря
1\,сут), возврат медленный (30\,сут): защёлка на восстановлении, а не на нагрузке.
\emph{Дофаминовая петля} (v5.1): тонус $d$, движимый растормаживанием $E/E_r-1$ или
долей $D$ ($\tau_d=7$\,сут), входит в нагрузку множителем $(1+a_d\,d)$; при $a_d=0$
модель воспроизводится побитово.

\subsection{Переключатель сна и бодрствования}\label{sec:sleep}
Филлипс--Робинсон, набор «new nominal» (Скелдон и соавт., 2014), взятый целиком, без
настройки:
\begin{equation}\label{eq:pr}
\begin{aligned}
&Q(V)=\frac{Q_{\max}}{1+e^{-(V-\theta)/\sigma}},\quad
\tau_v\dot V_v=\nu_{vm}Q(V_m)+D_v-V_v,\quad
\tau_m\dot V_m=A_m+\nu_{mv}Q(V_v)-V_m,\\
&D_v=B_0+\nu_{vc}\cos\tfrac{2\pi(t-\alpha)}{24}+A_G G-k_d D_{\mathrm{cpl}},\quad
\dot G=\begin{cases}(1-G)/18{,}2\,\mathrm{ч}&\text{бодрствование},\\-G/4{,}2\,\mathrm{ч}&\text{сон},\end{cases}
\end{aligned}
\end{equation}
бодрствование $:=\;Q(V_m)>1$\,с$^{-1}$. Сверка с публикацией: складки на
$1{,}450/2{,}463$\,мВ (в статье $1{,}45/2{,}46$); свободный цикл
$15{,}5/8{,}5/24{,}0$\,ч; размах сомногенов $0{,}527$ против $0{,}50$ --- без
подгонки. Повреждение входит членом $-k_d D_{\mathrm{cpl}}$ ($k_d=2$\,мВ): повреждённая
популяция подавляет сонное ядро и удлиняет бодрствование.

\subsection{Кора}\label{sec:cortex}
Пара Вильсона--Коуэна ($E$ --- возбуждение, $I$ --- PV),
$F(z,a,\theta)=1/(1+e^{-a(z-\theta)})$:
\begin{equation}\label{eq:wc}
\dot E=-E+F(c_1E-c_2gI+P_e),\qquad
\dot I=-I+F(c_3E-c_4gI),\qquad
g=1-0{,}4\,D,
\end{equation}
$P_e=P_{e0}+P_{e1}\min(Q(V_m)/Q_{\mathrm{wake}},1)$. Множитель на \emph{все} выходные
веса PV с $g_{\mathrm{dep}}=0{,}4$ взят из Вольман и соавт., 2011 (торможение до
60\,\%). Времена коры в $10^3$ раз быстрее всего остального: стационар затухающей
итерацией (адиабатическое исключение); гамма-ритма нет намеренно. При полном
повреждении $E$ вырастает на 19\,\% --- следствие двух литературных чисел, а не
подобранное усиление.

\subsection{Популяция}\label{sec:pop}
$J=32$ клетки, отличающиеся одним скаляром --- множителем нагрузки
$v_j=\exp(\sigma_v z_{p_j})/\mathrm{mean}$, $\sigma_v=0{,}10$, закреплённая по краям
сетка вероятностей, привязанная к $J=32$. Клетки связаны только через активность сети
(нагрузка на уцелевших растёт с долей повреждённых). Плавность --- от разброса, а не от
сети (Ф5, опровергнуто). Тонкость сетки: $J$ --- параметр модели через реализованный
разброс, $\mathrm{std}(\ln v)=0{,}1226/0{,}1060/0{,}0980/0{,}0943$ при
$J=8/16/32/64$ против $0{,}0910$ в пределе (сходимость $O(1/J)$), при чувствительности
по логарифму $\approx2{,}5$ (раздел~\ref{sec:J}).

\subsection{Нагрузка, эпизоды и закон сохранения разрывов}\label{sec:gap}
Нагрузка $L(t)$: исходный уровень на $q=0{,}10$ окна; суточный размах 20\,\% (вес $A_H$
калибруется по переключателю); быстрый шум $\sigma=0{,}08$ по логарифму нагрузки,
время корреляции $0{,}3T$; медленный шум \emph{выключен} (обоснование снято в Ф6,
Ф18); необязательные серотониновые всплески ($\times1{,}5$, частота $0{,}1$/сут, спад
2\,ч, потолок $\times2$); вспомогательные протоколы для правил о петлях (двухфазный
«сломать --- отпустить»; постоянный множитель $n_{\mathrm{mul}}$ как статический
эквивалент петли).

\begin{definition}[Эпизод, полоса]
Эпизод --- время, пока доля повреждённых $D>1/2$ (для сравнений по $J$ --- непрерывная
переменная среднего повреждения; момент пересечения порога интерполируется внутри
шага). \emph{Полоса} --- множество наборов, у которых разрыв вверх перекрывается
доступным всплеском, а разрыв вниз укладывается в бюджет шума $2\sigma+\sigma_s$.
\end{definition}

\begin{proposition}[Сохранение разрывов]\label{prop:gap}
При отношении границ окна $W$, положении нагрузки $q$ (логарифмическая шкала),
множителе обратной связи $F$, суточных пике и впадине $P,T$:
\begin{equation}\label{eq:gap}
\ln\frac{W^{1-q}}{P}+\ln(W^q F T)=\ln\frac{WFT}{P}\quad(\text{от }q\text{ не зависит}).
\end{equation}
\end{proposition}
У типичного набора сумма $\approx9\sigma$ шума нагрузки; клиническая картина (редкие
эпизоды длиной часы--сутки) требует $\approx4{,}5\sigma$. \emph{Один и тот же шум не
может одновременно редко срывать и быстро отпускать.} Выходы: односторонний триггер
либо переменная защита, поднимающая собственную нижнюю границу, --- выбран второй
(раздел~\ref{sec:defence}).

\section{Вычислительная часть}\label{sec:code}
\subsection{Устройство кода}
\begin{center}\footnotesize
\begin{tabular}{@{}>{\raggedright\arraybackslash}p{0.30\textwidth}>{\raggedright\arraybackslash}p{0.66\textwidth}@{}}
\toprule
Файл & Роль\\
\midrule
\code{run.py} & точка входа: разбор ключей, самопроверка критериев, блоки, отчёт, журнал\\
\code{config/base.toml} & все параметры прогона ($n_{\mathrm{sel}}=2400$, 60\,сут,
$M=1600$, $J=32$, разгон 16\,сут, $\sigma=0{,}08$, $k_d=2{,}0$, $\sigma_v=0{,}10$,
$\tau_s=2{,}625$\,сут)\\
\code{px/consts.py} & постоянные и сигмоиды, источник каждого числа в комментарии\\
\code{px/core.py} & кривая стационаров, окно, выборка наборов, сетка уязвимости\\
\code{px/sleep.py} & переключатель: складки, свободный цикл, калибровка\\
\code{px/cortex.py} & стационар Вильсона--Коуэна\\
\code{px/model.py} & \code{simulate}: один шаг по всем наборам и блокам перебора\\
\code{px/fast.py} & слитое numba-ядро ($\times3{,}3$) и быстрый \code{simulate}\\
\code{px/conf.py} & чтение TOML, кэш наборов, слои, закон разрывов, полоса\\
\code{px/batch.py} & блоки, склейка по хешу, распараллеливание по блокам, кэш блоков
(ключ --- хеш исходников модели)\\
\code{px/stats.py} & критерии: знаковый, граница медианы, CV$^2$ промежутков,
вынесение вердикта с проверкой MDE, самопроверки\\
\code{px/predict.py} & аналитические ожидания до всякого счёта (модель риска)\\
\code{px/report.py} & отчёт прогона и журнал вердиктов \code{ledger.json}\\
\code{px/rules/fNN.py} & по правилу на файл: \code{needs(cfg)}, \code{run(ctx)}\\
\code{px/rules/_blocks.py}, \code{_meas.py} & общие конструкторы блоков; сводка по блоку\\
\code{mass_core6.py} & замороженный монолит-эталон, в счёте не участвует\\
\code{tests/test_regression.py}, \code{tests/test_fast.py} & побитовые сверки с
монолитом и ядром\\
\bottomrule
\end{tabular}
\end{center}

\subsection{Конвейер, блоки, правила}
\begin{verbatim}
python run.py --list                 # какие правила есть
python run.py --rules all --dry      # оценка блоков и времени, без счёта
python run.py --rules f19            # одно правило; f02,f06,f18: несколько
python run.py --ledger               # накопленный журнал вердиктов
python tests/test_regression.py      # побитовая сверка с mass_core6.py
\end{verbatim}
Считаются только запрошенные блоки; одинаковые блоки склеиваются по хешу (78
запрошенных $\to$ 40 уникальных). Результаты блоков кэшируются под ключом с хешем
исходного кода модели: правка модели обесценивает кэш по построению. Распараллеливание
--- только по блокам: шум разыгрывается на длину одного блока и повторяется по точкам
перебора, поэтому блоки строго парные. Полный прогон на 2400 наборах $\approx3$\,ч
(слитое ядро). Дисциплина правил: (i) блоки --- через общие конструкторы в
\code{_blocks.py}: один и тот же блок для всех правил, иначе парные сравнения
развалятся; (ii) каждая часть --- через \code{stats.adjudicate} или
\code{stats.sign\_adjudicate}, которые возвращают «непроверяемо» \emph{симметрично},
когда минимально различимый эффект больше порога; (iii) каждая часть пишет
\code{ctx.verdict(...)} в \code{ledger.json} (дата, хеши блоков) --- вердикты не
пересматриваются; дефект отмечается рядом с вердиктом, вопрос перезадаётся новым
правилом на новом сиде.

\subsection{Статистика}\label{sec:stats}
Критерий пригоден, только если до счёта известны его нулевое значение, его шум при
нашем числе наборов и его мощность:
\begin{equation}\label{eq:mde}
\mathrm{MDE}=(z_{0{,}95}+z_{0{,}80})\,\mathrm{SE}=2{,}4865\,\mathrm{SE};
\end{equation}
утверждения о согласии используют допуск $3\,\mathrm{SE}$, обязанный быть шире MDE,
иначе утверждение неопровержимо в принципе. Знаковый критерий несёт \emph{обязательное
направление} (однажды молчаливая ошибка знака превратила $p=0{,}038$ в $p=0{,}989$);
медианы --- по односторонней границе порядковой статистики; промежутки --- по CV$^2$
внутри набора (нуль ровно единица); каждый критерий самопроверяется на синтетике с
известным ответом до всякого прогона. Аналитические ожидания (\code{px/predict.py}):
логарифм нагрузки стоит на $m_0=\ln(\mathrm{lo}\cdot(\mathrm{hi}/\mathrm{lo})^q)$,
барьер --- на $\ln(\mathrm{hi})$; срыв есть выброс шума за барьер, частота
$h(t)=e^{-d(t)^2/2}$ при $d(t)=(G-\ell(t))/\sigma$ (предэкспонента сокращается во всех
отношениях). Из одной этой модели, без подгонки: доля срывов во второй половине
бодрствования (Ф2а: $0{,}82$), кривая отклика на размах, вилка середины шкалы (Ф5),
доля срывов от всплесков (Ф8а), удлинение бодрствования прямо из детерминированного
переключателя (Ф4а, Ф6а), аффинный закон длительности (Ф7, Ф8б). Ранговая корреляция по
восьми ячейкам имеет ошибку $1/\sqrt{7}=0{,}38$; выход --- наклоны по наборам с
ослаблением $\rho_{\mathrm{наб}}=\rho_{\mathrm{свод}}/\sqrt{1+\lambda}$: в настоящих
данных $\lambda=33{,}8$, MDE падает с $0{,}94$ до $0{,}076$, и четыре части из пяти
закрываются.

\subsection{Воспроизводимость}\label{sec:repro}
(i) \emph{Замороженный монолит}: \code{tests/test_regression.py} требует побитового
совпадения выборки наборов, арифметики полосы и \code{simulate} с \code{mass_core6.py}
на закреплённых конфигурациях; сверка между машинами --- по допуску $10^{-12}$ с
отпечатком машины (измеренное наихудшее $3\cdot10^{-15}$). (ii) \emph{Слитое ядро}:
собрано из исходника эталона подстановкой текста с проверкой единственности каждого
куска (правка модели ломает сборку, а не расходится молча); сверено по всем 67 полям
на девяти закреплённых конфигурациях (счётные побитово; наихудшее расхождение
непрерывных $9{,}5\cdot10^{-15}$ при допуске $10^{-12}$) и по диффу всех вердиктов;
выигрыш в $3{,}3$ раза по настенным часам. (iii) \emph{Шаг по времени}: непрерывные
средние сходятся до $0{,}5$\,\% уже при $M=800$; счётные --- нет: предел Ричардсона на
$2{,}0$\,\% выше значения при $M=1600$, поэтому порогам по счётным величинам тоньше
$\approx3$\,\% верить нельзя --- это записано, а не замазано.

\subsection{Семь статистических ловушек}\label{sec:traps}
Каждая однажды дала ложный вывод в наших руках; каждая закрыта проверкой в коде:
(1) медиана по наборам от доли редких событий вырождена; (2) доля коротких промежутков
тождественно равна $1+\mathrm{CV}^2$ частот наборов --- физики не видит; (3) направление
знакового критерия --- обязательный аргумент; (4) нулевая ошибка по построению в
столбцах «в долях собственного уровня»; (5) насыщающаяся измеряемая величина топит
настоящий эффект --- мерить частоту, а не долю повреждавшихся; (6) допуск
равносильности шире собственного размаха кривой неопровержим; (7) среднее по
бимодальной популяции ($D>0{,}9$ против $D<0{,}01$) --- это доля сломанных наборов, а
не состояние.

\section{Результаты}\label{sec:results}
\subsection{Журнал вердиктов}
128 предзарегистрированных частей в 26 правилах (Ф2--Ф29): 66 выполнены, 48 не
выполнены, 13 непроверяемы, одна ниже собственной аналитической вилки. В таблице ---
доли частей: В --- выполнено, НВ --- не выполнено, НП --- непроверяемо, Ви --- ниже
вилки.

\begin{center}\footnotesize
\begin{tabular}{@{}l p{0.52\textwidth} c@{}}
\toprule
Правило & Вопрос & В/НВ/НП/Ви\\
\midrule
Ф2 & срывы копятся к концу бодрствования; порог по размаху & 2/4/0/0\\
Ф3 & повреждение откладывает сон, удлиняет эпизоды & 8/0/1/0\\
Ф4 & сон портится до порога срыва (предвестник) & 3/3/1/0\\
Ф5 & плавность доли повреждённых & 0/1/0/1\\
Ф6 & предвестник: высота, наклон, медленный шум & 4/2/2/0\\
Ф7 & длительность в пределах клинических двух суток & 3/3/0/0\\
Ф8 & серотониновый всплеск как триггер & 2/7/0/0\\
Ф10 & состав эпизода: хвост против времени над верхом окна & 3/2/0/0\\
Ф11 & индукция закрывает эпизоды & 1/2/0/0\\
Ф12 & глухой (рефрактерный) период от защиты & 1/2/0/0\\
Ф13 & независимость от числа клеток $J$ & 2/2/0/0\\
Ф14 & память защиты маскируется шумом & 2/1/0/0\\
Ф15 & глухой период есть, и он не от защиты & 3/0/0/0\\
Ф16 & аффинный закон длительности, проверка вне подгонки & 1/2/0/0\\
Ф18 & кто делает предвестник & 2/2/2/0\\
Ф19 & прайминг доводит до клинической цели & 4/0/0/0\\
Ф20 & ритм ломается лестницей захвата & 7/0/0/0\\
Ф21 & третья медленная величина --- заполнение популяции & 2/1/0/0\\
Ф22 & $J$ через эффективный разброс (проверка подменой) & 1/1/0/0\\
Ф23 & асимметричная защита следует $w(r)$ & 3/0/1/0\\
Ф24 & рваная нагрузка бьёт сильнее гладкой & 0/2/2/0\\
Ф25 & три фазы; форма полосы & 3/2/0/0\\
Ф26 & дофаминовая петля против статического эквивалента & 2/2/0/0\\
Ф27 & окно сосуществования нагрузочной петли & 2/3/0/0\\
Ф28 & петля ёмкости, первая редакция (контроли не прошли) & 0/3/4/0\\
Ф29 & петля ёмкости, только сопоставленные различия & 5/1/0/0\\
\bottomrule
\end{tabular}
\end{center}

\subsection{Что устояло}\label{sec:survived}
\begin{enumerate}[leftmargin=1.2em]
\item \emph{Длительность задаёт память защиты, а не возможности починки.} 99{,}8\,\%
дисперсии длительности внутри набора объясняется ходом нагрузки, 0{,}2\,\% ---
починкой (разложение точное по построению, доли дают в сумме единицу с точностью до
$10^{-16}$). Закон
аффинный:
\begin{equation}\label{eq:dur}
\text{длительность}=1{,}28+0{,}97\,\tau_b\ \text{(суток)},
\end{equation}
подгонка ниже 100\,ч, проверка в точке 150\,ч вне подгонки; насыщающаяся форма с
общим сдвигом отвергнута (Ф16б), хотя количественное предсказание вне подгонки
промахнулось мимо своей полосы $2\,\mathrm{SE}$ (Ф16а). Пассивная починка
промахивается мимо двух суток в $2$--$10$ раз (нужная память $\approx12$\,ч меньше
первой ступени).
\item \emph{Медленная память защиты делает предвестник:} при её отнятии высота падает
у 107 из 174 наборов, $p=0{,}0015$; медленный шум нагрузки не нужен.
\item \emph{Три фазы, а не две.} Полоса целых клеток, \emph{уязвимее} своих рухнувших
соседей: $0{,}00732\pm0{,}00042$ популяции ($17\sigma$; $\approx8{,}5$\,\% на мелкой
сетке); ровный нуль при замороженной защите; в $1{,}82$ раза шире при прайминге.
Гормезис на уровне популяции, предсказание модели, а не артефакт счёта.
\item \emph{Прайминг (Ф19) доводит до клинической цели:} медиана $1{,}46$\,сут при
односторонней 95-процентной границе $1{,}665$; коротких эпизодов $0{,}363$; парный
$p=5{,}2\cdot10^{-6}$; все четыре части прошли, повторено на 2000 наборах. Без
подгонки: потолок $\times3$ измерен, $\tau_p$ перебиралась; ускорение первой ступени
даёт от $-1$ до $-5$\,\% против $-52$\,\% у усиления, как и предсказано. Механизм:
прайминг поднимает \emph{верх} окна, а не ускоряет починку.
\item \emph{Асимметричная защита (Ф23):} перебор $r=0{,}33$, $0{,}5$, $1$, $2$, $3$
предсказывает $w=0{,}441$, $0{,}544$, $0{,}705$, $0{,}827$, $0{,}878$
по~\eqref{eq:wr}; измеренные уровни $0{,}0575$, $0{,}0644$, $0{,}0792$, $0{,}0958$,
$0{,}1065$ --- порядок точный ($\rho=+1$,
перестановочный $p=0{,}0083$). Доля повреждённых падает $0{,}486\to0{,}457$
($p=0{,}048$); длительность $7{,}74\to6{,}88$\,сут. Равенство постоянных подъёма и
спада было \emph{выбором режима}, а не упрощением; $r$ обязан остаться параметром
модели, пока измерение не задаст его значение.
\item \emph{Петля ёмкости (Ф29) превращает частые короткие эпизоды в редкие длинные:}
длительность $13{,}1\to25{,}6$\,сут, а частота срывов \emph{падает}
$0{,}042\to0{,}033$ на набор-сутки под риском; эпизодичность $0{,}75$ от разомкнутой
модели; сломанное состояние удерживается после сброса нагрузки, монотонно по глубине
сброса ($\rho=+1{,}000$, $p=0{,}0083$); прогрессирование превосходит сопоставленную по
ёмкости статику на $+0{,}0231$ при MDE $0{,}0096$, воспроизведено как $+0{,}0223$ и
$+0{,}0211$ на двух новых сидах. Единственная не прошедшая часть: \emph{форма} риска
во времени не меняется (тренд риска $1{,}94$ против порога $2{,}23$, Ф29в).
\item \emph{Ритм: лестница захвата} с рациональными числами вращения и чередованием
длин циклов на дробных ступенях (раздел~\ref{sec:staircase}).
\item \emph{$J$ входит только как эффективный разброс стойкости;} предел $J\to\infty$
существует с поправками $1/J$ (раздел~\ref{sec:J}).
\end{enumerate}

\subsection{Что опровергнуто}\label{sec:failed}
\begin{itemize}
\item \emph{Нагрузочная петля типа Грейс (Ф26, Ф27):} необратимость включается при
силе петли $1{,}11$, а эпизоды гаснут при $0{,}18$ --- шестикратный структурный
разрыв, устойчивый к более мягкому протоколу; окна, где хроничность и эпизодичность
сосуществуют, нет. Разомкнутая модель отпускает все сломанные наборы
($0{,}9936\to0{,}0000$), петля отпускает $0{,}24$\,\%. Причина: нагрузка одна на
набор, петля движет популяцию целиком.
\item \emph{Петля ёмкости:} уровень и длительность меняются, форма риска во времени
--- нет (выше).
\item \emph{Рваная нагрузка не бьёт сильнее гладкой (Ф24):} частота срывов падает на
40\,\% с ростом времени корреляции при закреплённых среднем и дисперсии; доля
повреждавшихся насыщается.
\item \emph{Серотониновый всплеск (Ф8):} за всплеском следует 20\,\% срывов против
требуемых 80\,\% --- ручка громкости, а не спусковой крючок; ни одно свойство эпизода
не меняется.
\item \emph{Суточный размах (Ф2б):} плавный рост в $1{,}5$ раза, ворот нет; плавность
--- от разброса, а не от сети (Ф5, наблюдение \emph{ниже} вилки).
\item \emph{Эпизод не помнит глубину (Ф10а):} длительность совпадает со временем над верхом
окна (ранговая корреляция $0{,}94$ на прежней модели / $0{,}87$ на нынешней); хвост
$\approx7$\,\% эпизода.
\item \emph{Форма полосы} $\varphi(\Phi^{-1}(1-D))$ опровергнута дважды с
противоположными знаками (Ф25г,д) --- ширина не является функцией $D$.
\item \emph{Проверка подменой сигмы (Ф22б) не прошла:} эффективный разброс объясняет
измеренное падение, но не закрывает вопрос (кандидаты: дисперсия среднего множителя
$1/J$ и дискретность $D$ в обратной связи на сон).
\item Несколько частей о наклонах предвестника непроверяемы при выбранном разрешении
сетки (MDE больше записанного порога).
\end{itemize}

\subsection{Отдельные механизмы}\label{sec:selected}
\paragraph{Лестница (Ф20).}\label{sec:staircase}
Детерминированный переключатель под 24-часовой силой, повреждение тормозит сонное
ядро: число вращения стоит на рациональных значениях и перескакивает:
\begin{equation}\label{eq:stair}
\begin{gathered}
1{:}1\to5{:}4\to3{:}2\to2{:}1\to3{:}1\to4{:}1\to\dots\to9{:}1\\
\to\text{сна нет},
\end{gathered}
\end{equation}
переходы при $kD\approx4{,}77$; $4{,}78$; $4{,}79$; $5{,}09$; $5{,}16$\,мВ; пять
ступеней (5:1--9:1) внутри $5{,}172$--$5{,}178$\,мВ. Прогон с мелким перебором: 52
ступени в $0{,}9$\,мВ, у 34 чередуются длины циклов (узоры $25,24,24,47$\,ч при $4{,}772$;
$47,25$ при $4{,}779$; $46,26,24$ при $4{,}783$); ширины $0{,}3$\,мВ (2:1),
$8\cdot10^{-3}$ (3:2), $5\cdot10^{-6}$ (11:8). Лестница захвата (языки Арнольда), а
\emph{не} каскад удвоений: период прибавляет по пропущенной ночи; различающий признак
--- захват на периоде, не равном степени двойки (найден). Ограничение калибровочное:
тождественно $kD\le k$, поэтому при $k=2$\,мВ первая ступень ($4{,}77$\,мВ)
недостижима ни при каком повреждении; при $k=5$\,мВ модель до неё доходит
($kD=4{,}93$, $D=0{,}999$).

\paragraph{Число клеток $J$ (Ф13, Ф22).}\label{sec:J}
Три пути: (1) хвост распределения уязвимости --- закрыт закреплённой сеткой (при
$J=32$ совпадает с прежней до последнего знака); (2) квантование порога эпизода
($0{,}625$ при $J=8$ против $0{,}516$ при $J=64$) --- обойдено непрерывной
переменной; (3) реализованный разброс (раздел~\ref{sec:pop}): 11\,\% по эффективной
$\sigma_v$ между $J=16$ и $64$ при чувствительности $\approx2{,}5$ $\Rightarrow$
$\approx25$\,\% по измеряемым величинам --- ровно тот размах, который Ф13 нашла на
2000 наборах. Дискретизация может точно сохранять среднее и края и всё равно
промахиваться по дисперсии ($O(1/J)$): сравнивать прогоны с разным $J$ --- значит
сравнивать разные модели.

\paragraph{Петля ёмкости (Ф28, Ф29).} Первая редакция провалилась на собственных
контролях: «при разомкнутой петле прибавка нулевая» --- утверждение ложно: бистабильный
набор не отпускается, доля сломанных растёт сама; контроль прошёл на малой выборке и
провалился
по мере накопления данных. Поэтому в Ф29 нет ни одного критерия против нуля: все
утверждения --- сопоставленные различия на одном уровне повреждения
(раздел~\ref{sec:lessons}).

\paragraph{Полоса (Ф25).} На сетке в 600 узлов: замороженная защита $\Rightarrow$
точная монотонность (ноль нарушений); живая защита рвёт порядок на $11{,}9$\,\% шагов,
худшее нарушение $0{,}973$ (полный переворот); прайминг усиливает, но не создаёт
($A=0$: $0{,}0184$). Следствие: замкнутая формула доли повреждённых $1-F_v(v^*)$ при
живой защите незаконна; любая редукция обязана нести явный пограничный слой из трёх
конфигураций.

\section{Методологические уроки}\label{sec:lessons}
В модели с поглощающими состояниями и разнородной популяцией \emph{абсолютного нуля
для накопительных величин не существует}: доля повреждённых растёт сама (бистабильный
набор не отпускается), риск падает сам (уязвимые поглощаются первыми). Два наших
контроля утверждали, что эти базовые уровни равны нулю; оба были ложны, и оба \emph{прошли на
малой выборке и провалились по мере накопления данных}: ложный нуль проходит по
нехватке мощности --- порог сжимается как $1/\sqrt{n}$, а базовый уровень стоит на
месте. Всякое
утверждение обязано быть сопоставленным различием плеч, приведённых к одному уровню
повреждения. Семь ловушек (раздел~\ref{sec:traps}) --- тот же урок в семи обличьях; ещё
одна: размах $n$ зашумлённых значений сам случаен
($\mathbb{E}[\text{размах}]=1{,}693\,r$ при $n=3$), поэтому порог «разброс ниже
10\,\%» был бессмыслен при $r=12$\,\% и стал проверяем только на 2000 наборах --- где
немедленно провалился и вскрыл третий путь $J$.

\section{Область применимости и открытые вопросы}\label{sec:scope}
У модели нет стационара без поглощения: она применима в окнах, коротких по сравнению
со временем накопления, --- это модель \emph{возникновения и окончания эпизода}, а не
течения болезни. Петля ёмкости --- единственный проверенный механизм, дающий
прогрессирование, и делает он это удлинением эпизодов, а не их учащением. Открыто:
(i) прогрессирование в масштабе болезни требует звена, которое не возвращается
полностью, --- намеренно не добавлено: механизм, добавленный ради нужного результата,
--- не проверка; сначала пишется правило, которое его опровергнет;
(ii) области, связанные через узкие места, --- против однородности «нагрузка одна на
набор», убившей нагрузочную петлю; (iii) поломка ритма недостижима при $k=2$\,мВ
против первой ступени $4{,}77$\,мВ --- неверная калибровка или другой путь;
(iv) ширину полосы задаёт что-то третье (вероятнее всего, время у самого
опрокидывания) --- проверяемо против уже измеряемого времени в середине шкалы,
обязано совпасть на обоих переборах, иначе снимается; (v) третья медленная величина
предвестника (у заполнения популяции нет своего $\tau$); (vi) свободный
член~\eqref{eq:dur} ($1{,}3$\,сут при нулевой памяти) не объясняется $x$;
(vii) кинетика пероксиредоксина измерена на клеточных линиях, а не на PV-клетках;
(viii) дофамин и симптомы --- вне модели по замыслу.

\section{Воспроизводимость}\label{sec:run}
Команды и конфигурация: раздел~\ref{sec:code}. Вердикты: \code{ledger.json} с датами
и хешами блоков. Эталон: \code{mass_core6.py} (заморожен); побитовые сверки:
\code{tests/}. Полный прогон на 2400 наборах $\approx3$\,ч; всегда сначала
\code{--dry}. Хеш исходного кода модели входит в ключ кэша блоков: правка модели
обесценивает кэш по построению.

\smallskip
\noindent{\small\emph{Благодарности.} Текст и код моделирования создавались с помощью
Claude Opus 5 и DeepSeek, как и в сопроводительной статье \code{paper.tex}. Все числа
взяты из журнала вердиктов, отчётов и комментариев к исходному коду проекта;
биологические опорные факты снабжены метками происхождения в брошюре
\code{brochure.tex}.}

\smallskip
\begin{thebibliography}{9}\small
\itemsep0pt
\bibitem{hardingham} G.~E. Hardingham, K.~Q. Do, \emph{Nat. Rev. Neurosci.} 17, 125 (2016).
\bibitem{cuenod} M. Cuenod et al., \emph{Mol. Psychiatry} (2021--2022).
\bibitem{selvaggio} G. Selvaggio, P. Coelho, A. Salvador, \emph{Redox Biol.} (2018).
\bibitem{volman} V. Volman, M.~M. Behrens, T.~J. Sejnowski, \emph{J. Neurosci.} (2011).
\bibitem{skeldon} A.~C. Skeldon, D.-J. Dijk, G. Derks, \emph{PLoS One} 9, e103877 (2014).
\bibitem{daan} S. Daan, D.~G.~M. Beersma, A.~A. Borb\'ely, \emph{Am. J. Physiol.} 246, R161 (1984).
\bibitem{howes} O.~D. Howes et al., \emph{Arch. Gen. Psychiatry} (2012).
\bibitem{papadia} S. Papadia et al., \emph{Nat. Neurosci.} 11, 476 (2008).
\bibitem{kil} I.~S. Kil et al., \emph{Mol. Cell} (2015).
\bibitem{delolmo} M. del Olmo et al., \emph{Int. J. Mol. Sci.} (2019).
\bibitem{harkness} J.~H. Harkness et al., \emph{bioRxiv} (2020), publ.\ (2021).
\bibitem{bischoff} Bischoff et al., raw sulfiredoxin induction kinetics, \emph{Sci. Rep.}.
\bibitem{koulakov} A.~A. Koulakov et al., \emph{Nat. Neurosci.} 5, 775 (2002).
\end{thebibliography}

\end{document}
